\documentclass[../main.tex]{subfiles}

\begin{document}

\subsection{Memory Boundedness, Basic Tiling, and Analysis}

\begin{frame}
  \frametitle{Let's do some experimenting}
  \textbf{This kernel assigns one thread per output element, and computes them
  independently. So changing block size shouldn't matter right?}

  \centering
  \includegraphics[width=0.72\textwidth]{generated/naive_variant_latency.png}
\end{frame}

\begin{frame}
  \frametitle{Blocks and SMs}
  \begin{block}{Interpretation}
    This happens because entire \textbf{blocks} are scheduled to SM's. The larger
    block sizes of 16x16 and 32x32 lead to more \textbf{L1 cache hits}.
  \end{block}

  \pause
  \begin{table}
    \centering
    \begin{tabular}{c|c}
      \textbf{Block shape} & \textbf{L1 load hit rate} \\
      \hline
      $8\times8$ & 83.51\% \\
      $16\times16$ & 87.13\% \\
      $32\times32$ & 94.98\%
    \end{tabular}
  \end{table}
  \centering
  {\small RTX 3090, $M=N=K=4096$}
\end{frame}

\begin{frame}
  \frametitle{Shared Memory (SMEM) Caching}
  \begin{block}{Programmable L1 Cache (SMEM)}
    Remember how our launch config required us to group threads as blocks?
    That is because \textbf{blocks} are scheduled to SM's as a unit, so \textbf{threads
    in the same block have access to the SM's L1 cache/SMEM.}
  \end{block}

  \centering
  \includegraphics[width=0.90\linewidth]{cuda-programming-model.png}

  {\tiny Image source: \href{https://modal.com/gpu-glossary/device-software/cuda-programming-model}{Modal, ``CUDA Programming Model,'' \emph{GPU Glossary}.}}
\end{frame}

\begin{frame}
  \frametitle{Tiling}
  Explicit caching of data shared between threads.
  \begin{block}{Main Idea}
    Before doing any computation, load a chunk of the necessary data from DRAM
    to SMEM as a block. Then, memory accesses for the computation now only go to SMEM.
  \end{block}
\end{frame}

\begin{frame}
  \frametitle{Tiling steps}

  \centering
  \only<1>{
    \includegraphics[height=.7\textheight]{tiled_steps/0_load_smem.png}
  }
  \only<2>{
    \includegraphics[height=.7\textheight]{tiled_steps/0_partial_sum.png}
  }
  \only<3>{
    \includegraphics[height=.7\textheight]{tiled_steps/0_advance_tile.png}
  }
  \only<4>{
    \includegraphics[height=.7\textheight]{tiled_steps/1_load_smem.png}
  }
  \only<5>{
    \includegraphics[height=.7\textheight]{tiled_steps/1_partial_sum.png}
  }
  \only<6>{
    \includegraphics[height=.7\textheight]{tiled_steps/2_repeat_0.png}
  }
  \only<7>{
    \includegraphics[height=.7\textheight]{tiled_steps/2_repeat_1.png}
  }

\end{frame}

\begin{frame}[fragile]
  \frametitle{Basic tiled SGEMM}

  \begin{columns}[T,onlytextwidth]
    \begin{column}{0.75\textwidth}
      \begin{lstlisting}[
        basicstyle=\ttfamily\fontsize{5.4}{6.1}\selectfont,
        escapeinside={(*@}{@*)}
      ]
__shared__ float As[BK][BK];
__shared__ float Bs[BK][BK];
float sum = 0.0;

for (int k = 0; k < K; k += BK) { (*@\tikz[remember picture,overlay]\coordinate (tile-loop);@*)
  const int A_tile_x = k + xt;
  if (y < M && A_tile_x < K) {
    As[yt][xt] = A[y * K + A_tile_x];
  } else {
    As[yt][xt] = 0.0;
  }

  const int B_tile_y = k + yt;
  if (B_tile_y < K && x < N) {
    Bs[yt][xt] = B[B_tile_y * N + x];
  } else {
    Bs[yt][xt] = 0.0;
  } (*@\tikz[remember picture,overlay]\coordinate (tile-load);@*)

  __syncthreads();

  for (int i = 0; i < BK; i++) {
    sum += As[yt][i] * Bs[i][xt]; (*@\tikz[remember picture,overlay]\coordinate (tile-compute);@*)
  }
  __syncthreads();
}
      \end{lstlisting}
    \end{column}

    \begin{column}{0.25\textwidth}
      \begin{tikzpicture}[remember picture,
          callout/.style={draw,rounded corners,fill=#1!10,
          text width=0.92\linewidth,inner sep=3pt,font=\scriptsize}
        ]
        \node[callout=codecomment] (loop-note) {\textbf{Outer loop:}\par
        Loop over tiles.};
        \node[callout=codecuda,below=5mm of loop-note] (load-note) {
        \textbf{Cooperative loading to SMEM}};
        \node[callout=orange,below=5mm of load-note] (compute-note) {
        \textbf{Compute partial sum only with SMEM}};
      \end{tikzpicture}
    \end{column}
  \end{columns}

  \begin{tikzpicture}[remember picture,overlay,
    pointer/.style={-{Stealth[length=2mm]},thick}]
    \draw[pointer,codecomment] (loop-note.west) -- (tile-loop);
    \draw[pointer,codecuda] (load-note.west) -- (tile-load);
    \draw[pointer,orange!85!black] (compute-note.west) -- (tile-compute);
  \end{tikzpicture}

  {\tiny Source: \texttt{my\_sgemm/src/sgemm.cu}.}
\end{frame}

\begin{frame}
  \frametitle{Basic tiled SGEMM visualization}
  \begin{itemize}
    \item Watch each block load small regions of $A$ and $B$ together.
    \item The loaded tiles are reused to update an entire tile of $C$.
    \item Synchronization separates loading, computing, and advancing to the next tiles.
  \end{itemize}

  \vfill
  \centering
  \href{https://drive.google.com/file/d/10VVCEAk6f6WpjGJXf4Ojl5NytiW02ofl/view?usp=sharing}{\beamerbutton{Play tiled-kernel visualization}}
  \vfill
\end{frame}

\begin{frame}{Tiling cuts latency by another order of magnitude}
  \centering
  \includegraphics[width=0.7\textwidth]{generated/cpu_naive_tiled_latency.png}
  22.6\% improvement in latency at \(n,m,k=8192\).
\end{frame}

\begin{frame}
  \frametitle{Can we still do better?}

  \begin{block}{Throughput analysis}
    We report a final latency of \(0.36\)s or \(360\)ms for \(n,m,k=8192\).
    This is around \(10\times\) the theoretical minimum of $31$ms. This slowdown
    suggests we are still \textbf{memory bound}.
  \end{block}

  \pause
  This tiling setup with \(32 \times 32\) blocks and \(n, m, k = 8192\), we need
  a \textbf{grid} of \(\frac{m}{32} \times \frac{n}{32} = 256\times 256 = 65,536\) blocks.

  \pause
  Each block loads \(\frac{k}{32} = 256\) tiles from both \(A\) and \(B\).
  Each tile is \(32\times 32\) elements.

  \pause
  Thus, we have
  \[
    \underbrace{65,536}_{\text{blocks}}
    \times \underbrace{256}_{\text{tiles/block}}
    \times \underbrace{2}_{A\text{ and }B}
    \times \underbrace{32 \times 32}_{\text{elements/tile}}
    \times \underbrace{4\text{ bytes}}_{\text{bytes/element}}
    \approx 137.7\text{ GB}.
  \]

  \pause
  At \(935.8 \text{ GB/s}\), we need 0.147s or 147ms to move all data from DRAM
  to SMEM. \pause \textbf{40\% of our latency is moving data.}
\end{frame}

% Not necessary.
% \begin{frame}
%   \frametitle{This type of analysis is getting hard...}
%   When we did this analysis for the \textbf{naive kernel}, it was easy to see how
%   the data moved, and at the end of it, we even got the nice result of
%   \[
%     \underbrace{3.55\text{s}}_{\text{Observed latency}}
%     \approx \underbrace{3.52\text{s}}_{\text{Time to move data}}
%     + \underbrace{0.031\text{s}}_{\text{Theoretical minimum for work}}.
%   \]

%   \pause However, as we apply more optimizations, \textbf{predicting performance
%   based on hardware specifications becomes impractical.} \pause Even with just
%   our \textbf{basic tiling kernel}, the observed latency of
%   \[
%     \underbrace{360\text{ms}}_{\text{Observed latency}}
%     \ne \underbrace{147\text{ms}}_{\text{Time to move data}}
%     + \underbrace{32\text{ms}}_{\text{Theoretical minimum for work}}
%     = 179\text{ms}.
%   \]

%   \pause This gap comes from things like occupancy, synchronization overhead, bank conflicts...
%   \pause \textbf{Without profiling every possible source of latency, can we still
%   determine whether memory is the bottleneck}?
% \end{frame}

\end{document}
